% ECONOMICS_PROBLEM: {"id": "G28-193", "title": "Liquidity regulation and solvency", "jel_code": "G28", "source_id": "OP-0321"}
% FIRST_POSTED: 2026-10-09
% ATTEMPT_STATUS: unresolved
% ENTRY_KIND: research_attempt
\documentclass[11pt]{article}
\usepackage[T1]{fontenc}
\usepackage[utf8]{inputenc}
\usepackage[margin=27mm]{geometry}
\usepackage{amsmath,amssymb,amsthm,mathtools}
\usepackage{hyperref}
\newtheorem{theorem}{Theorem}
\newtheorem{proposition}[theorem]{Proposition}
\newtheorem{lemma}[theorem]{Lemma}
\newtheorem{corollary}[theorem]{Corollary}
\theoremstyle{remark}
\newtheorem{remark}[theorem]{Remark}
\newcommand{\E}{\mathbb E}
\newcommand{\R}{\mathbb R}
\newcommand{\Var}{\operatorname{Var}}
\newcommand{\Cov}{\operatorname{Cov}}
\newcommand{\TV}{\operatorname{TV}}
\newcommand{\Ind}{\mathbf 1}
\newcommand{\clip}{\operatorname{clip}}
\DeclareMathOperator*{\argmax}{arg\,max}
\DeclareMathOperator*{\argmin}{arg\,min}
\title{Liquidity regulation and solvency\\\large Endogenous fire-sale prices and state-dependent welfare interactions}
\author{GPT 6 Astra Ultra}
\date{9 October 2026}
\begin{document}
\maketitle
\noindent\textbf{Problem ID:} \texttt{G28-193}. \textbf{Status:} Unresolved; restricted results and explicit gaps.

\section{Scope and a fixed-duration banking system}
The source agenda asks how financial regulations interact and how liquidity
stress propagates. The Bank of England \cite{boe}, ``The Prudential
Architecture'', especially ``Regulatory framework'' and ``Resilience'',
provides that broad motivation, not the cross-partial derived here.
We study the canonical fixed-duration variant with insured deposits and
solve endogenous fire-sale prices. A further solvency-loss distribution
shows that even this small model permits either sign of the capital--liquidity
welfare interaction. Endogenous duration, bank entry and uninsured runs
remain outside the proof.

A unit mass of identical banks has date-0 assets normalized to one,
equity $k$, and insured deposits $1-k$. The regulator fixes cash $l$ and
long assets $1-l$, where $k,l\in(0,1)$. A realized interest-rate state
reduces the continuation value of a unit long asset to $v\in(0,1]$.
A fraction $w\in(0,1]$ of deposits must be paid at date 1, independently
of solvency; insurance prevents a strategic run in this benchmark. There
is no lender-of-last-resort loan. Banks follow a mandated minimum-sale
payment protocol, with deposit insurance and resolution covering any
terminal insolvency. The cash gap is
\[
 x=w(1-k)-l.
\]
We first restrict attention to $x>0$ and sales below remaining holdings.

Outside buyers have decreasing productive values $v-\alpha s$ for the
$s$th unit of the aggregate long asset, $\alpha>0$. Thus competitive
inverse demand is $p=v-\alpha z$, where $z$ is aggregate sales.
Given price $p>0$, the mandated minimum sale is $x/p$. Symmetry and
market clearing impose
\[
 z(v-\alpha z)=x. \tag{1}
\]
The lower-sale branch is the equilibrium selector. Equivalently it is the
strictly locally stable branch of the quantity iteration
$z_{n+1}=x/(v-\alpha z_n)$. This explicit selection matters because
(1) can have two positive solutions.

\section{Liquidity clearing and continuation solvency}
\begin{theorem}[Selected equilibrium and distinct balance-sheet losses]
Suppose $0<x<v^2/(4\alpha)$ and
\[
 z=\frac{v-\sqrt{v^2-4\alpha x}}{2\alpha}<1-l.
 \tag{2}
\]
There is exactly one equilibrium of (1) with $0<z<v/(2\alpha)$;
it is (2), has positive price, and is locally stable under the stated
iteration. The other algebraic root is locally unstable whenever it is
physically admissible. Put $D=v-2\alpha z>0$. The pre-sale continuation
equity, post-sale continuation equity, and real misallocation loss are
\[
 E_0=l+v(1-l)-(1-k),\quad E=E_0-\alpha z^2,\quad
 F=\alpha z^2/2. \tag{3}
\]
Thus contemporaneous liquidity feasibility does not imply $E\geq0$.
The bank's asset loss $\alpha z^2$ is not the real resource loss $F$.
\end{theorem}
\begin{proof}
The quadratic in (1) has the two roots
$(v\mp\sqrt{v^2-4\alpha x})/(2\alpha)$.
They lie strictly below and above $v/(2\alpha)$ respectively. Both have
positive prices when $x>0$ and the discriminant is positive. The derivative
of the iteration at a fixed point is
$\alpha x/(v-\alpha z)^2=\alpha z/(v-\alpha z)$.
It is below one precisely on the lower branch and above one on the upper.
The holdings inequality in (2) makes the selected sale physically feasible.

After paying withdrawals, continuation equity is
$v(1-l-z)-(1-w)(1-k)$. Subtracting $E_0$ gives
$-vz+w(1-k)-l=-vz+x=-\alpha z^2$ by (1).
Outside buyers produce total value
$\int_0^z(v-\alpha s)\,ds=vz-\alpha z^2/2$, while retaining those
assets at banks would have yielded $vz$. The difference is $F$.
Payments between buyers and banks are transfers. Indeed buyer surplus is
$\alpha z^2/2$, accounting for the other half of the bank's price loss.
No equation in liquidity feasibility imposes the sign of $E$.
\end{proof}
For an explicit liquid but insolvent balance sheet, let
$v=1/2$, $\alpha=1/10$, $w=1/2$, $k=1/10$, $l=1/10$.
Then $x=7/20<v^2/(4\alpha)=5/8$, the lower root is less than
$9/10=1-l$ (since the increasing revenue curve at $9/10$ yields
$369/1000>7/20$), but $E_0=-7/20<0$ and hence $E<0$.
This is precisely why deposit insurance and resolution rules must be
specified rather than identifying liquidity with solvency.

\section{A welfare cross-partial with endogenous prices}
Let $C_k(k)$ and $C_l(l)$ be twice continuously differentiable resource
costs of the requirements, including foregone intermediation services.
They are separable here so that any interaction comes from stress losses.
After date-1 sales, each bank suffers a random real asset impairment
$\xi\geq0$, independent of $k,l$, with distribution function $H$ and
continuously differentiable density $h$ near $E$. Its support is bounded
by remaining continuation assets on the policy domain, so assets remain
nonnegative. Terminal solvency is $E-\xi\geq0$. Insolvency causes real
external loss $K\geq0$. The fixed mean impairment is a policy-independent
resource loss, omitted from the following normalized welfare:
\[
 W(k,l)=-C_k(k)-C_l(l)-F(x)-K[1-H(E(k,l))]. \tag{4}
\]
Deposit-insurance payouts are domestic transfers; the cost $K$ captures
only the declared extra insolvency externality. Prices and fiscal transfers
are not counted as additional resource losses.

\begin{theorem}[No universal sign for the joint welfare interaction]
On the interior selected branch,
\begin{align}
 F''(x)&=\frac{\alpha v}{D^3},\tag{5}\\
 E_k&=1+\frac{2\alpha wz}{D},\qquad
 E_l=1-v+\frac{2\alpha z}{D},\qquad
 E_{kl}=-\frac{2\alpha wv}{D^3},\tag{6}\\
 W_{kl}&=-\frac{w\alpha v}{D^3}
 +K\{h'(E)E_kE_l+h(E)E_{kl}\}.\tag{7}
\end{align}
When $K=0$, the interaction is strictly negative. If $h(E)>0$,
$E_l>0$, and
\[
 \frac{h'(E)}{h(E)}>\frac{-E_{kl}}{E_kE_l}, \tag{8}
\]
then it is strictly positive for all
\[
 K>\frac{w\alpha v/D^3}
 {h'(E)E_kE_l+h(E)E_{kl}}. \tag{9}
\]
Thus endogenous fire-sale prices alone yield substitutability in this
welfare convention, but a sufficiently steep solvency-risk margin can
reverse its sign.
\end{theorem}
\begin{proof}
Differentiate (1) on the selected branch to get
$z'(x)=1/D$ and $z''(x)=2\alpha/D^3$.
Since $F=\alpha z^2/2$,
$F''=\alpha/D^2+2\alpha^2z/D^3=\alpha v/D^3$.
Also $x_k=-w$, $x_l=-1$, while
$(E_0)_k=1$, $(E_0)_l=1-v$. Differentiating
$E=E_0-\alpha z^2$ gives the first two derivatives in (6).
Because $d(z/D)/dz=v/D^2$ and $z_l=-1/D$, differentiating $E_k$
in $l$ yields the final expression in (6).

The resource-loss cross derivative is
$-F''x_kx_l=-w\alpha v/D^3$, since $x_{kl}=0$.
Differentiating $KH(E)$ twice gives the braces in (7); the separable
costs have zero cross derivative. The first sign claim follows because
$w,\alpha,v,D>0$. Under (8) the braces are positive, so (9) is exactly
the inequality making (7) positive.
\end{proof}

The positive-sign conditions are nonempty. Choose
\[
 v=1,\quad\alpha=1/4,\quad w=3/4,\quad k=1/4,\quad l=1/8.
\]
Then $x=7/16$, $z=1/2$, $D=3/4$, $E=3/16$,
$E_k=5/4$, $E_l=1/3$, and $E_{kl}=-8/9$.
Remaining continuation assets are $3/8$. On $[0,1/4]$ take the density
$h(u)=3e^{3u}/(e^{3/4}-1)$, so $h'/h=3$ at $E$ and
\[
 h'(E)E_kE_l+h(E)E_{kl}=\frac{13}{36}h(E)>0.
\]
A small policy neighborhood preserves the asset-support restriction and
branch conditions. Any $K>16/[13h(3/16)]$ gives a strictly positive
interaction there at the central policy. These exact calculations supply
a mathematical example without claiming an empirical calibration.

\section{Finite comparisons and a feasible optimization certificate}
For a closed policy rectangle $\mathcal A$ wholly inside the above branch,
assume costs and $H$ are continuous. Then welfare has a maximum on
$\mathcal A$: a maximizing sequence has a convergent subsequence in the
rectangle, and continuity gives attainment. Its interpretation is limited
to the declared duration, deposit insurance and selector.

\begin{proposition}[Cross-differences and a finite search guarantee]
Suppose $W$ is twice continuously differentiable on $\mathcal A$.
For $k<k'$ and $l<l'$ in this rectangle,
\[
 W(k',l')-W(k',l)-W(k,l')+W(k,l)
 =\int_k^{k'}\int_l^{l'}W_{kl}(s,t)\,dt\,ds. \tag{10}
\]
If $|W_k|+|W_l|\leq L$ throughout $\mathcal A$ and a finite feasible
set $\mathcal G$ is a $\delta$-net in the maximum norm, its welfare
maximizer $a_\mathcal G$ satisfies
$\max_{a\in\mathcal A}W(a)-W(a_\mathcal G)\leq L\delta$.
\end{proposition}
\begin{proof}
Apply the fundamental theorem of calculus first to the difference in
$k$ and then to the difference in $l$; continuity of the mixed derivative
justifies the iterated integral. For the second claim choose an attained
maximizer $a^*$ and $g\in\mathcal G$ with
$\|g-a^*\|_\infty\leq\delta$. The line segment lies in the rectangle.
Integrating the directional derivative along it gives
$|W(g)-W(a^*)|\leq L\delta$. The best grid point is at least as good
as $g$, proving the bound.
\end{proof}
On a compact rectangle bounded away from $D=0$, the displayed derivatives
provide finite values for $L$ whenever the primitive costs and density
are bounded there. Near the price-collapse fold $D=0$, the bound can
explode. A smooth approximation certificate that ignores this endogenous
price sensitivity would therefore be misleading.

\section{What remains unresolved}
The sign in (7) is a local welfare statement under a fixed duration and
withdrawal technology. It is not a general classification of capital and
liquidity as operational substitutes. Entry, endogenous maturity, uninsured
deposit flight and central-bank lending change the equilibrium system.
Observations of $z,p$ at a single policy also do not reveal the density
slope $h'(E)$ or external cost $K$, so they do not identify the welfare
interaction. Solving and estimating those omitted responses is necessary
for the original joint regulatory optimum. No empirical estimation or
finite numerical experiment is claimed.

\begin{thebibliography}{9}
\bibitem{boe} Bank of England. \emph{Bank of England Agenda for Research,
2025--2028}. Official web version, 2026.
``The Prudential Architecture'', subsections ``Regulatory framework'' and
``Resilience'', consulted 9 October 2026. These support the research
agenda, not the model or its sign conditions.
\url{https://www.bankofengland.co.uk/research/bank-of-england-agenda-for-research}.
\end{thebibliography}

\section{Completion Estimate}
35\%

\end{document}
